\section{Implementation}

% Explain how the prototype realises the design described in Chapter 4.
% Focus on implementation choices, component responsibilities, and how
% components interact. Do not repeat the protocol rationale or report results
% here; refer to the System Design chapter and present evaluation evidence in
% the Results chapter.

\subsection{Prototype Overview}
\textbf{The prototype is built using the following tools and technologies:}
\begin{itemize}
    \item \textbf{Git} for version control and repository-based storage and exchange;
    \item \textbf{Python} for the client-side protocol operations;
    \item \textbf{HTML, CSS and JavaScript} for the feed and directory websites;
    \item \textbf{GitHub} for hosting user and directory repositories, with
    GitHub Pages serving the static websites; and
    \item \textbf{PyNaCl and GitPython} for cryptographic and Git operations
    in the Python client.
\end{itemize}

\textbf{The repository contains the \texttt{client} source code,
\texttt{thesis\_files} for the thesis, and documentation such as
\texttt{README.md}, which explains how to run the prototype, and
\texttt{latest\_design.md}, which contains the workflow diagrams. The
\texttt{docs/archive} directory retains earlier design documents. The client
entry point is \texttt{client/app.py}. Running \texttt{python app.py} from the
\texttt{client} directory displays the available commands, including profile
creation, social actions, replication, feed display and site publication. The
automated tests use \texttt{pytest}; their results are reported in the Results
chapter.}

% Keep this overview concise: the System Design chapter already explains the
% protocol architecture and workflows. This chapter describes how the
% prototype realises them.

\subsection{Identity and Signed Data}

\textbf{The \texttt{profile create} command is handled by
\texttt{ProfileCreator}, which clones the canonical repository template and
generates an Ed25519 key pair using PyNaCl. It creates a profile containing
the user's handle, public key and profile information, then signs and saves
it as \texttt{social/profile.json}. The private key is stored separately in
the client's identity state, while the public key and signature are included
in the profile. The user then runs \texttt{publish} with the URL of the Git
repository where the identity will be hosted. \texttt{PublishManager}
publishes the repository, adds the repository URL to the profile, re-signs
and pushes the updated profile, and registers the account in the social
directory.}

\textbf{The \texttt{Signer} and \texttt{Verifier} classes provide the shared
JSON signing and verification operations. Before signing, \texttt{Signer}
serialises the object with sorted keys and excludes any existing signature.
It signs these bytes with the Ed25519 private key and stores the resulting
signature in Base64 form. \texttt{Verifier} decodes the signature and checks
it against the same sorted JSON representation using the public key. The
following lines show the core operations (a screenshot of these lines can be
included here):}

\begin{lstlisting}[caption={Core JSON signing and verification operations}]
message = json.dumps(obj_to_sign, sort_keys=True).encode("utf-8")
signed = self.signing_key.sign(message)
self.verify_key.verify(message, signature_bytes)
\end{lstlisting}

\textbf{\texttt{ProfileVerifier} uses \texttt{Verifier} to check a profile's
signature before returning its verified handle, public key and repository
URL. This is the implementation of the profile verification step shown in
Figure~\ref{fig:identity_workflow}.}

\subsection{Social Action Implementation}
\textbf{Social actions are stored as JSON files in the user's
\texttt{social/actions/} directory. As described in the System Design
chapter, they share common fields: an identifier, action type, author public
key, creation timestamp and signature. Since these fields and the creation
steps are common across post, reply, like and follow actions, the
implementation uses \texttt{ActionBase} to avoid duplicating this logic.}

\textbf{As shown in Figure~\ref{fig:action_class_diagram},
\texttt{ActionBase} defines the shared creation workflow in
\texttt{\_create()}. It loads the public key from \texttt{profile.json} and
the private key from the client's identity state, checks that the keys
correspond, generates an action ID and builds the common fields. It then calls
\texttt{\_extend()}, allowing each subclass to add its action-specific
fields: \texttt{PostAction} adds post content; \texttt{ReplyAction} adds
content, the \texttt{inReplyTo} action ID and target handle;
\texttt{LikeAction} adds the target action ID and target handle; and
\texttt{FollowAction} adds the followed user's public key as its target.
Finally, \texttt{ActionBase} signs and writes the JSON file. The client uses
\texttt{GitPublisher} to commit and push the created action.}

\textbf{This shared creation workflow uses the Template Method pattern:
\texttt{ActionBase} defines the common algorithm, while subclasses customise
the \texttt{\_extend()} step \cite{gfg_template_method}. Separately, the
prototype uses a factory design pattern to avoid duplicating the logic for
instantiating different action classes. When a user enters a command,
\texttt{ActionFactory.create(...)} obtains the corresponding class from
\texttt{ActionRegistry} and creates an instance. This allows \texttt{app.py}
to request an action by command without knowing the concrete class or its
construction details, encapsulating object creation and decoupling command
handling from action-specific implementation. The registry maps command names
to action classes, avoiding a growing chain of action-specific
\texttt{if}/\texttt{elif} branches in the factory
\cite{gfg_factory_method}.}

\textbf{This structure supports extension as new social actions are introduced.
To add an action, a developer can implement a class inheriting from
\texttt{ActionBase} and register it for a command; the existing factory and
its dispatch logic do not need to change. The new command must still be
connected to the CLI parser and registry, but no new action-specific branch
needs to be added to the factory. This is an example of the Open--Closed
Principle: the action system can be extended without modifying its existing
creation logic. The Post, Reply, Like and Follow sequence diagrams in the
System Design chapter
(Figures~\ref{fig:post_workflow}, \ref{fig:reply_workflow},
\ref{fig:like_workflow} and \ref{fig:follow_workflow}) show how these action
types are used in the protocol workflows.}

\begin{figure}[h]
    \centering
    \includegraphics[width=\textwidth]{figures/action_class_diagram.png}
    \caption{Social Action Class Structure}
    \label{fig:action_class_diagram}
\end{figure}


\subsection{Following and Discovery}
\textbf{Following builds on the shared action-creation process described
above. \texttt{FollowAction} creates a signed Follow action whose target is
the followed user's public key. In addition, the follow command uses the
repository URL supplied by the user to configure a Git remote, allowing later
replication to fetch from that repository. The implementation follows the
validation and follow sequence shown in Figures~\ref{fig:discovery_workflow}
and \ref{fig:follow_workflow}; the key implementation steps are described
below.}

\textbf{Before accessing the target repository, the client retrieves the
handle's entry from the public directory and checks that the supplied
repository URL matches the registered URL. The directory response is obtained
over HTTP using \texttt{urllib.request}; the configured URL identifies the
published JSON data. The client then checks that the supplied repository URL
matches the entry's \texttt{repoURL}, rejecting a mismatch before fetching the
repository.}

\begin{lstlisting}[caption={Checking a Follow target against the directory}]
with urllib.request.urlopen(DIRECTORY_JSON_URL, timeout=5) as response:
    directory = json.load(response)
entry = directory.get(handle)
if not isinstance(entry, dict):
    raise ValueError(f"Handle '{handle}' is not registered")
if entry.get("repoURL") != repo_url:
    raise ValueError("Repository URL does not match the directory")
\end{lstlisting}

\textbf{This excerpt shows the two checks in order: the client downloads the
published JSON, selects the entry indexed by the requested handle, then
rejects a missing entry or a repository URL that differs from the registered
one. Only after these checks does it access the target repository.}

\textbf{After this check, the client adds the target repository as a named Git
remote, using the handle as the remote name. A Git remote is a saved
name-to-URL association in the local repository; it does not copy the remote
actions at this point. GitPython's \texttt{create\_remote()} adds this
configuration, which the replication process can later use to fetch the
remote branch. This reuses Git's existing remote mechanism rather than
maintaining a separate list of followed repository URLs
\cite{github_remotes}.}

\begin{lstlisting}[caption={Adding a followed repository as a Git remote}]
if handle not in [r.name for r in repo.remotes]:
    repo.create_remote(handle, repo_url)
\end{lstlisting}

\textbf{The condition avoids adding the same named remote twice; when absent,
\texttt{create\_remote()} stores the handle-to-URL association in the local
repository's Git configuration. The remote name is subsequently used to
fetch that repository.}

\textbf{The client then fetches the target's \texttt{main} branch and reads
\texttt{social/profile.json} from it. \texttt{ProfileVerifier} delegates
signature checking to the shared \texttt{Verifier}; the follow operation
also compares the verified handle and public key, and the profile's repository
URL, with the requested values and the directory entry before creating the
Follow action. Thus a valid signature proves control of the profile's key,
while the directory comparison checks that this is the key registered for
the handle.}

\begin{lstlisting}[caption={Validating the fetched Follow profile}]
verified = ProfileVerifier(profile).verify()
if verified["handle"] != handle:
    raise ValueError("Signed profile handle does not match")
if verified["publicKey"] != registered_identity["publicKey"]:
    raise ValueError("Profile key does not match the directory")
if profile["repoURL"] != repo_url:
    raise ValueError("Profile repository URL does not match")
\end{lstlisting}

\textbf{Here, the first line validates the profile's signature. The following
comparisons then bind that signed profile to the requested handle, the
directory's registered public key and the repository URL supplied by the
user. A valid signature alone establishes control of a key; these additional
comparisons establish that the key and repository correspond to the handle
being followed.}

\textbf{Once the profile passes these checks, \texttt{FollowAction} creates
the signed action targeted at the verified public key. The action is
committed and pushed to the follower's repository. Fetching the remote
\texttt{main} branch to obtain the profile can also transfer Git objects from
its history, but Follow does not import the remote actions into the local
action set used by the client. The saved remote is used by the separate
replication process to discover, verify and accept actions for local use.
This keeps Follow responsible for establishing the relationship and
replication responsible for processing remote activity, as described in the
corresponding System Design workflows.}

\textbf{The directory supports this lookup by mapping each handle to a
repository URL and registered public key. When an account is published,
\texttt{PublishManager} checks an existing entry before updating it. If the
handle already has a registered key, publication refuses to replace it with
a different key. This is the publishing-side safeguard that complements the
Follow-time comparison and helps detect a recreated repository presenting a
different key for the same handle.}

\begin{lstlisting}[caption={Preventing replacement of a registered key}]
registered_key = entry.get("publicKey")
if registered_key != public_key:
    raise RuntimeError(
        f"Refusing to replace the registered public key for '{handle}'. "
        "Use the identity recovery process to change keys."
    )
\end{lstlisting}

\textbf{This guard runs when the account is published: if the handle already
has a directory entry, the stored key must match the key being published.
The exception prevents an ordinary publish operation from silently replacing
the handle's established key.}

\textbf{The directory entry also contains an optional \texttt{siteURL}.
After site publication succeeds, \texttt{PublishManager} adds this URL to
the existing entry, allowing the directory website to link the handle to
the published feed. The repository URL remains available separately for
following and protocol access.}

\textbf{The directory website is a static interface to this JSON data.
\texttt{index.html} defines elements with the IDs \texttt{member-list},
\texttt{member-count} and \texttt{directory-status}. In \texttt{app.js},
\texttt{document.querySelector()} uses these IDs to select the elements
where the directory cards, profile count and loading or empty-directory
message are displayed. The loader fetches and sorts entries from
\texttt{directory.json}, then renders a card for each handle. The following
excerpt shows the connection between the HTML elements and the main display
updates:}

\begin{lstlisting}[caption={Selecting and updating directory page elements}]
const memberList = document.querySelector("#member-list");
const directoryStatus = document.querySelector("#directory-status");
const memberCount = document.querySelector("#member-count");
memberList.replaceChildren(...entries.map(
  ([handle, entry]) => renderMember(handle, entry)
));
memberCount.textContent = `${entries.length} ${entries.length === 1 ? "profile" : "profiles"}`;
directoryStatus.textContent = entries.length === 0
  ? "No profiles are listed in the directory yet."
  : "";
\end{lstlisting}

\textbf{\texttt{replaceChildren()} fills the member-list element with the
cards created by \texttt{renderMember()}, while the next lines update the
count and show a message if there are no entries. Within each card, a valid
\texttt{siteURL} links the handle to the published feed and the repository
URL is shown as a separate link. \texttt{style.css} controls the visual
layout. The website makes directory entries discoverable; the client's
Follow validation—not the web page—checks that a handle, key and repository
agree.}

% Optional: include a screenshot of the rendered social-directory website
% (rather than a source-code repository screenshot) to show how users access
% repository and published-feed links.
% \begin{figure}[h]
%     \centering
%     \includegraphics[width=\textwidth]{figures/social_directory.png}
%     \caption{Static social-directory website}
%     \label{fig:implementation_social_directory}
% \end{figure}

\subsection{Replication}

% Describe how Replicator obtains remote repository data, verifies candidate
% actions, stores accepted actions locally, and propagates them. Explain the
% role of ActionVerifier and profile verification where relevant. Refer to
% Figures~\ref{fig:replication_workflow} and
% \ref{fig:multihop_replication} in the System Design chapter.

\subsection{Feed Construction and Site Publication}

% Describe how locally available signed actions and verified profile
% information are assembled for display, and how the static site is published
% separately from protocol data. Identify the responsibilities of
% ShowFeedAction, SitePublisher, and PublishManager accurately; distinguish
% local feed display from generation/publication of feed.json. Refer to
% Figure~\ref{fig:feed_construction_workflow} in the System Design chapter.

\subsection{Implementation Verification}

% Briefly describe the test strategy, test environment, and how tests map to
% important components or requirements. Report actual test outcomes, scenario
% evidence, and requirement assessments in the Results chapter rather than
% predicting or asserting them here.

\subsection{Implementation Limitations}

% Describe implementation-specific constraints that affect interpretation or
% reproducibility. Reserve broader implications and future work for Discussion.